We study whether markets and later innovation outcomes distinguish between more credible and less credible AI disclosure in annual 10-K filings. Using U.S. public firms from 2016 to 2025, we classify AI-related sentences as actionable, speculative, or irrelevant, aggregate those classifications into firm-year disclosure measures, and use them to construct the AS ratio and PatentMismatch. Results show that low-credibility AI disclosure becomes common quickly in the post-2022 period: PatentMismatch characterizes 41.5\% of AI-talking firm-years in 2024 and 42.1\% in 2025. That pattern is difficult to reconcile with a view in which the recent expansion of AI language in corporate filings simply mirrors a comparable expansion in verifiable technological implementation.

The market evidence indicates that this divergence is not fully resolved at the filing date. Filing-date abnormal returns show little clean evidence of immediate discipline, whereas the stronger correction appears after disclosure. A long--short strategy that buys firms with more credible AI disclosure and shorts mismatch firms earns its strongest return over the subsequent quarter, with an equal-weight alpha of 0.93\% per month, or roughly 11.2\% annualized. The corresponding value-weight result is much weaker, Suggesting that the slower correction is concentrated outside the largest firms. Cross-sectional valuation evidence points in the same direction. At the same time, the disclosure construct remains anchored to later technological realization: actionable AI language aligns more closely with subsequent AI patenting than speculative language, and the AS-ratio $\times$ PatentMismatch specification predicts weaker future AI patent output. Exploratory post-ChatGPT tests are directionally consistent with this interpretation, although they are better understood as suggestive rather than as the primary identification strategy.

The contribution lies in the combination of these results with the underlying measurement design. Prior work has shown that AI-related language in corporate disclosure is informative and that markets react to AI narratives, while recent studies have begun to distinguish AI talk from AI walk more directly. The present paper contributes by centering mandatory annual filings, scaling the actionable-versus-speculative distinction into a reusable audited framework, and introducing PatentMismatch as a filing-based construct that combines disclosure composition with contemporaneous patent weakness. This makes it possible to study low-credibility AI disclosure in a way that is portable across samples and informative both for later innovation outcomes and for market consequences.

Several limitations remain. Because PatentMismatch firms may differ systematically from non-mismatch firms along observable and unobservable dimensions, the return-based evidence is best interpreted as consistent with slower incorporation of disclosure credibility rather than as definitive proof that the pricing differential is caused solely by AI-washing. The post-ChatGPT tests should likewise be understood as exploratory rather than decisive. Some of the market-based patterns are also materially stronger outside the largest firms, which suggests that mispricing and disclosure credibility may be most relevant in settings where valuation is harder and external verification is weaker. The evidence suggests that low-credibility AI disclosure is economically important not only because it departs from later technological realization, but also because markets do not appear to fully price that gap immediately. More broadly, the results point to a setting in which the cost of AI talk can fall more quickly than the cost of real AI implementation, and in which disclosure composition provides a useful way to distinguish the two.